% ─────────────────────────────────────────────────────────────────────────────
\section{Individual Reflections}
\label{sec:reflect-learning}

\subsection{Muhammad Ahmad Humayun Hanif}

I came into this project thinking the hard part would be the AI. It turned out the hard part was the database. Once we decided that a slot document would serve as both the inventory record and the booking record, everything else had to be built around that decision. The OCR flow, the agent state machine, the vendor calendar, they all read and write that same document. When any of them fell out of step, it was immediately visible and genuinely painful to fix. That taught me more about system design than anything I have read, because I felt the consequences directly.

LangGraph took longer to get right than I expected. The concept of a persistent graph that resumes across multiple webhook calls sounds clean on paper, but debugging a cyclic state machine where a node can be re-entered several calls later is a different kind of problem. I learned to add explicit logging at every edge transition just to know where the graph was at any given moment. I would not skip that step again.

\subsection{Jazib Waqas}

The backend work on the booking pipeline was the most demanding part of the project for me. Getting the concurrency control to work correctly took longer than I initially planned. The theory was straightforward but the implementation had a lot of edge cases, particularly around making sure that a transaction retrying automatically would not create a duplicate record. Once that was solid, everything built on top of it was much cleaner.

The thing I would change if we started over is agreeing on a proper API contract between backend and frontend from day one. We did not do that early on and it cost us real time mid-project when the frontend was calling endpoints that had changed. We eventually sorted it out, but we should have treated the schema as a shared document that neither side could change without the other knowing.

\subsection{Taha Hunaid Ali}

The frontend was harder than I expected, mostly because the data requirements were more complex than a typical app. Slot availability can change at any moment from three different sources simultaneously: a WhatsApp booking, an app booking, and a walk-in. Keeping the UI consistent without hammering the backend or confusing the user required thinking about caching more carefully than I had on previous projects. I rewrote the data layer twice before I was happy with it.

The social features were honestly the most fun to build and they got the best reaction from people we showed the app to. Getting the leaderboard to load quickly on mobile took real work though. The first version was loading everything at once and it was slow enough that it was embarrassing. Pagination fixed it, but I should have thought about that from the start.

\subsection{Muhammad Taqi}

My personal learning experience throughout this final year project was an incredible deep dive into the practical realities of Conversational AI, which is an area my experience has been growing in the past year. I primarily worked on the NLP model and LangGraph orchestration, moving from raw data annotation to deploying a functional state machine. Annotating and cleaning the data across English and Roman Urdu was tedious, but it was my first real lesson in how the quality of data absolutely defines the ceiling of a model's performance.

Writing the normalization layers and watching the LangGraph orchestrator successfully manage conversational state, hold context, and route users across multiple webhook calls was a highly rewarding progression. Debugging complex state leakage and cyclic memory bugs pushed my problem-solving skills, ultimately teaching me that the true reliability of an AI agent lives entirely in its deterministic state management rather than just the generative model wrapping it.

% ─────────────────────────────────────────────────────────────────────────────
\section{Team Reflection}
\label{sec:reflect-team}

We split the project by ownership, where each person was responsible for their component end-to-end. That worked well for moving in parallel but it created integration points that needed deliberate attention. We learned to treat those handoffs more carefully as the project progressed, but there were a few weeks early on where the backend and frontend were drifting apart without either side realising it.

The most stressful moment was the Twilio sandbox number expiring close to an evaluation. It was a one-person fix but the timing was bad. After that we put API credentials into a shared checklist that we reviewed before any major milestone.

The two biggest decisions we made that were not in the original proposal were switching from PostgreSQL to Firestore and deciding to fine-tune a small model rather than use a large one end-to-end. Both were driven by practical constraints that became obvious only after we started building. Firestore's transaction model was a better fit for the slot locking logic than anything we would have written on top of a relational database. And using Groq for everything would have been too slow and too expensive to run on free-tier infrastructure. In both cases the constraint pushed us to a better solution than the one we had originally planned.

If we did this again, we would write the API contract first and treat it as a binding document. That single change would have saved the most time.

%%% Local Variables:
%%% mode: latex
%%% TeX-master: "../report"
%%% End:
